\documentclass[a4paper,11pt]{article}

\usepackage{revy}
\usepackage[utf8]{inputenc}
\usepackage[T1]{fontenc}
\usepackage[danish]{babel}


\revyname{MatematikRevyen}
\revyyear{2025}
% HUSK AT OPDATERE VERSIONSNUMMER
\version{0.1}
\eta{$5$ minutter}
%\status{Færdig}

\title{Moderne bevisførelse}
\author{Carl '21}

\begin{document}
\maketitle

\begin{roles}
\role{I}[Snow] Instruktør
\role{F}[Johan R] Forelæser
\role{C}[Daniel] Censor
\role{S1}[Alberte F] Studerende
\role{S2}[Oliver] Studerende
\role{S3}[Elmer] Studerende
\end{roles}

\begin{props}
\item Tavle
\item Bord med to stole
\item Stopur
\end{props}

\begin{sketch}
\scene{Mundtlig eksamen. L og C sidder ved et bord. Der ligger to bunker
med spørgsmål på bordet. S1 kommer ind og giver dem håndenS}
\says{L} Velkommen til eksamen til An0. Som du ved, har KU’s ledelse besluttet, at alt for mange studerende består deres mundtlige eksamener
her på matematisk institut. Det er uklart hvilke tal de baserer deres vurdering på, men lad det nu ligge. Du kan bare træde nærmere og trække et spørsmål og et dertilhørende benspænd.
\scene{S1 trækker et spørgsmål fra den ene bunke og viser det}
\says{C} Spørgsmål 3... Det er Hovedsætningerne! Og hvordan skal det fremføres?

\scene{C gestikulerer mod den anden bunke og S1 trækker også et kort derfra. S1 læser kortet og ser forvirret ud}
\says{S1}[Tøvende] Bevis ved debat?? Hvad betyder det?

\says{L} UHHH! Ej, den er sjov!
\says{L}[Kigger begejstret på C] Går du ud og henter den næste person i rækken?

\scene{C går ud og henter S2 og stiller dem ved siden af S1}

\says{C} Så kan i lige slå sten, saks, papir om hvem, der er for og imod sætningen.

\scene{S1 og S2 slår om det. S1 vinder}

\says{L} Wuuh! Tillykke! Hvilken side vælger du så?

\says{S1}[Tøvende] ... for?

\says{S1} Men altså. Jeg synes, det virker en smule unfair det her...
\says{L}[afbryder] Shhhh! Det er slet slet ikke afgjort endnu. S1 du har 20 sekunder til dit åbningsargument. Begyndt
\scene{L tænder stopuret. S1 går lidt nervøst igang med at tegne billedet fra beviset for H2 i An0 bogen}

\says{S1}[langsomt] Hvis vi har givet en kontinuert funktion på et interval $[a,b]$ ind i R og et tal $\gamma$ mellem $f(a)$ og $f(b)$. Så ønsker vi at finde et $\xi$ så $f(\xi)=\gamma$ ...

\says{L} STOP! Det var tiiid. S2, din tur
\says{S2} Okay så. Tossen derovre prøver at overbevise jer om, at
jeg ikke kan gå herfra og hertil [peger på tavlen] uden at ramme $\gamma$. Men S1 har har overset, at jeg bare kan gøre sådan her!
\scene{S2 laver en kurve fra $f(a)$ ’uden om’ $\gamma$ ved at tegne en kurve der går venstre om y-aksen. L og C ser meget begejstrede ud og peger ivrigt på tavlen.}
\says{S1} Nej nej nej, det kan du ikke mene. For hvis jeg tager midtpunktet her mellem $a$ og $b$....
\scene{S1 kigger forvirret på tavlen og tvivler lidt. Men kommer så i tanke}
\says{S1} Og Ruselemmaet!! Hvis vi bruger ruselemmaet, så finder vi at...

\says{S2}[afbryder hånende] Ruse smruse. Hvis du har brug for et lemma, skal du så ikke lige beivse det først? Og jeg hører dig slet ikke argumenterer mod det her.

\scene{S2 peger på tavlen. L og C nikker til hinanden og virker enige med S2. S1 ved slet ikke, hvad de skal gøre.}
\says{S1} Men jeg... men jeg... fandt et midtpunkt...
\says{L} Okay, jeg tror, vi stopper den her. C og jeg er enige om, at vi har en klar vinder af den her debat. S2, tillykke. S1, vi ses til reeksamen.

\scene{S1 og S2 forlader lokalet. L og C sætter sig tilbage til bordet.}

\says{L} Næste!
\scene{S3 kommer ind, hilser og trækker spørgsmål. C kigger på spørgsmålet}
\says{C} Mmh, spørgsmål 6. Det er Taylor i flere variable. Og hvad er dit benspænd?
\scene{S3 trækker et kort og kigger med L på kortet. L ser begejstret ud, S3 ser forvirret ud.}
\says{S3} Greens sætning? 
\says{L} Du går bare igang. 

\scene{Tæppe}
\end{sketch}
\end{document}